\documentclass[11pt]{article}
\usepackage[a4paper,margin=24mm]{geometry}
\usepackage{amsmath,amssymb,amsthm,mathtools}
\usepackage[hidelinks]{hyperref}
\newtheorem{theorem}{Theorem}
\newcommand{\K}{K}
\title{An optimal lower bound for weighted Bessel tails}
\author{Research proof and Lean certificate; independent review pending}
\date{October 28, 2026}
\begin{document}
\maketitle
\begin{abstract}
For real $n<1$ and $\nu\ge 1-n+\sqrt{2(1-n)}$, we determine exactly which
orders $r$ make $x^{-\nu}K_r(x)$ a strict lower bound for
$\int_x^\infty t^{-\nu}K_{\nu+n}(t)\,dt$ at every positive $x$.
The answer is $|r|\le\sqrt{(\nu+n)^2-2\nu-1}$, including the zero-order
boundary. This proves Conjecture~2.4 of Gaunt's arXiv version~2.
The argument uses convexity in squared order and a localized integral
comparison for optimality. The accompanying Lean certificate proves both
directions directly from the classical integral representation of $K$.
\end{abstract}

\section{The statement}
Throughout, $K_\rho$ is the modified Bessel function of the second kind.
We use its classical representation \cite[Eq.~10.32.9]{DLMF}
\begin{equation}\label{eq:K}
 K_\rho(x)=\int_0^\infty e^{-x\cosh u}\cosh(\rho u)\,du,
 \qquad x>0,\quad \rho\in\mathbb R.
\end{equation}
Gaunt proposes the lower bound and the optimal order in
\cite[Conjecture~2.4 and Remark~2.5]{Gaunt}.
\begin{theorem}\label{thm:main}
Let $n<1$ and $\nu\ge1-n+\sqrt{2(1-n)}$, and put
$\mu=\nu+n$ and $\delta=\sqrt{\mu^2-2\nu-1}$.
For every $r\in\mathbb R$,
\begin{equation}\label{eq:iff}
 \left[\ \int_x^\infty t^{-\nu}K_\mu(t)\,dt
       >x^{-\nu}K_r(x)\ \text{for all }x>0\ \right]
       \quad\Longleftrightarrow\quad |r|\le\delta.
\end{equation}
If $|r|>\delta$, the reverse strict inequality holds for every sufficiently
large $x$.
\end{theorem}
The parameters are real; neither $n\ge0$ nor integrality is required.
The case $n=1$ is outside the theorem and is not claimed as a new result.

\section{The lower bound}
Fix $x>0$. The function $F(\lambda)=K_{\sqrt\lambda}(x)$ is strictly convex
on $[0,\infty)$. Indeed, for $u>0$,
\[
 \cosh(u\sqrt\lambda)=\sum_{j=0}^\infty
       \frac{u^{2j}\lambda^j}{(2j)!}
\]
is a sum of convex functions with a strictly convex quadratic term.
Integrating the strict inequality against the positive weight in
\eqref{eq:K} proves the claim. Convergence follows, for example, from
\[
 e^{-x\cosh u}\cosh(\rho u)
 \le \exp(|\rho|^2/x)\exp(-xu^2/4).
\]
The parameter assumptions give
\[
 0\le\delta<\nu,\qquad
 \mu^2=\delta^2+2\nu+1>(\delta+1)^2.
\]
To check the first assertion, use
$\delta^2=(\mu-1)^2-2(1-n)$ and
$\nu^2-\delta^2=(1-n)(2\nu+n+1)>0$.

For $\delta>0$, extrapolate the strictly convex function $F$ from
$(\delta-1)^2$ and $(\delta+1)^2$ to $\mu^2$. The result is
\begin{equation}\label{eq:chord}
 (\nu+\delta)K_{\delta+1}(x)+(\delta-\nu)K_{\delta-1}(x)
       <2\delta K_\mu(x).
\end{equation}
Integration by parts in \eqref{eq:K} gives
$K_{\delta+1}-K_{\delta-1}=2\delta K_\delta/x$.
Thus \eqref{eq:chord} becomes
\begin{equation}\label{eq:comparison}
 K_{\delta-1}(x)+\frac{\nu+\delta}{x}K_\delta(x)<K_\mu(x).
\end{equation}
There is no division by $\delta$ at the boundary. When $\delta=0$,
$\mu^2=2\nu+1$; the strict tangent inequality for
$\cosh(u\sqrt\lambda)$ at $\lambda=1$ gives
\[
 \cosh(\mu u)>\cosh u+\nu u\sinh u \qquad(u>0).
\]
Integration by parts for $u e^{-x\cosh u}$ yields
$\int_0^\infty u\sinh u\,e^{-x\cosh u}\,du=K_0(x)/x$.
After integration the same inequality \eqref{eq:comparison} follows, since
$K_{-1}=K_1$.

Differentiating \eqref{eq:K} under its integral, with domination at $x/2$,
gives $K_\delta'=-K_{\delta-1}-\delta K_\delta/x$.
Consequently \eqref{eq:comparison} implies
\[
 \frac{d}{dx}\bigl(x^{-\nu}K_\delta(x)\bigr)>-x^{-\nu}K_\mu(x).
\]
All terms are integrable on a positive half-line, and the boundary value
at infinity is zero. For example, for $t\ge a>0$,
$K_\rho(t)\le e^{-t/2}K_\rho(a/2)$.
Integrating proves the strict lower bound at $r=\delta$.
Evenness and monotonicity in absolute order, directly from \eqref{eq:K},
give it for every $|r|\le\delta$.

\section{Optimality by a local comparison}
Suppose $r>\delta$; evenness covers negative orders. Set
\[
 A=\frac{\nu+r}{2r},\qquad B=\frac{r-\nu}{2r},\qquad
 h(u)=\cosh(\mu u)-A\cosh((r+1)u)-B\cosh((r-1)u).
\]
Then $h(0)=h'(0)=0$ and $h''(0)=\delta^2-r^2<0$.
Two mean-value arguments give $\varepsilon>0$ such that
$h<0$ on $(0,\varepsilon]$. Choose $0<a<b<\varepsilon$ and $c>0$
with $h\le-c$ on $[a,b]$. The constant
$C=\int_0^\infty e^{-\cosh u}|h(u)|\,du$ is finite.
For $x\ge1$, positive integral comparison gives
\begin{equation}\label{eq:local}
 D(x):=\int_0^\infty e^{-x\cosh u}h(u)\,du
 \le-c(b-a)e^{-x\cosh b}+C e^{-(x-1)\cosh\varepsilon}.
\end{equation}
On $[a,b]$ use the negative constant bound; elsewhere below $\varepsilon$
the integrand is nonpositive. Above $\varepsilon$ its absolute value is
bounded by $e^{-(x-1)\cosh\varepsilon}e^{-\cosh u}|h(u)|$.
These observations also justify both tails in \eqref{eq:local}.
Since $\cosh b<\cosh\varepsilon$, we have $D(x)<0$ eventually.

The recurrence already proved identifies
\[
 D(x)=K_\mu(x)-K_{r-1}(x)-\frac{\nu+r}{x}K_r(x).
\]
For $x$ beyond that eventual-negative threshold, integrate $t^{-\nu}D(t)$
and use the weighted-$K$ derivative and zero boundary at infinity:
\[
 \int_x^\infty t^{-\nu}D(t)\,dt
   =\int_x^\infty t^{-\nu}K_\mu(t)\,dt-x^{-\nu}K_r(x)<0.
\]
This proves necessity in \eqref{eq:iff} and its eventual strict converse.

\section{Certificate and research status}
The Lean certificate uses \eqref{eq:K} as the definition of $K$ and proves
convergence, both integrations by parts, dominated argument differentiation,
tail integrability, strict comparisons, parameter geometry, and both
directions of \eqref{eq:iff}. Its endpoint is
\texttt{source\_full\_classification\_339406}. A fresh source replay checks
80 declarations in 14 modules with warnings as errors, using only
\texttt{propext}, \texttt{Classical.choice}, and \texttt{Quot.sound}.
No fixed-order asymptotic theorem is assumed in the certificate.

The source checked here is arXiv version~2; the journal-body comparison and
independent novelty/subsumption assessment remain open. In particular,
squared-order convexity and the other classical tools are not presented as
new machinery. The complete integrated certificate and this manuscript
await independent review. This draft has zero discovery credit, carries no
priority claim, and does not change any public problem status.

\begin{thebibliography}{9}
\bibitem{Gaunt} R. E. Gaunt, \emph{Inequalities for some integrals involving
modified Bessel functions}, Proc. Amer. Math. Soc. \textbf{147} (2019),
2937--2951, doi:\href{https://doi.org/10.1090/proc/14433}{10.1090/proc/14433}.
Statement checked in \href{https://arxiv.org/abs/1806.00524v2}{arXiv:1806.00524v2}.
\bibitem{DLMF} NIST Digital Library of Mathematical Functions,
\href{https://dlmf.nist.gov/10.32.E9}{Eq.~10.32.9}, classical integral
representation of $K_\nu$.
\end{thebibliography}
\end{document}
